\chapter{Tampon retained}
\index{Tampon retained}

\medskip
\noindent\textit{Educational reference; always seek professional advice for individual care.}

\section*{Definition and Overview}
A retained tampon, sometimes called a forgotten tampon, is a tampon that remains inside the vagina after its use, usually because the string has been lost, the tampon has been pushed out of reach, or the person simply forgot it was there. It is a surprisingly common problem, and it can happen to anyone who uses tampons, including experienced users. In most cases the tampon is harmless for a short time and can be removed easily, but the longer it stays in place, the greater the risk of infection and of a serious condition called toxic shock syndrome (TSS). Because a retained tampon can produce an unpleasant, sometimes blood-stained and smelly discharge, it is frequently mistaken for an infection such as thrush or bacterial vaginosis, and a tampon is then forgotten for even longer. The good news is that a retained tampon is simple to diagnose, simple to remove, and easily prevented with straightforward habits. This chapter explains what happens when a tampon is retained, the symptoms to look for, when to seek help, and how to prevent the problem altogether.

\section*{Epidemiology}
There are no national figures for retained tampons in many countries, because the problem is largely unreported, but it is a common presentation to general practitioners, sexual health clinics, and emergency departments. One United States survey found that a quarter of women reported having forgotten a tampon at some time. It is more likely to happen during heavy periods, when multiple tampons are changed quickly, or when a woman uses the highest absorbency products. The problem occurs most often towards the end of a period, when bleeding is light and a woman may not notice the absence of the tampon, or when she is busy, tired, travelling, or under stress. Women who use tampons for the first time, or who insert a tampon when bleeding has stopped, are also at risk. People of any age who use tampons, including teenagers, are susceptible, and adolescence is a common time for the first experience.

\section*{Aetiology and Pathophysiology}
Tampons are designed to be removed by pulling on a string that hangs outside the vagina. A tampon becomes retained when removal fails, for one of several reasons:

\begin{itemize}
    \item \textbf{The string is missing or lost:} the string may break, or the tampon may have been inserted without the string left accessible
    \item \textbf{The tampon was pushed too high:} a tampon pushed past the string, or insertion during intercourse, can lodge beyond easy reach
    \item \textbf{Simple forgetting:} the most common cause, particularly with light bleeding or when several days have passed
    \item \textbf{Insertion without an ongoing period:} a tampon inserted for discharge or at the very end of a period can be forgotten because there is no further bleeding to prompt its removal
\end{itemize}

Once retained, the tampon acts as a foreign body in a warm, moist environment rich in nutrients. Bacteria, particularly \textit{Staphylococcus aureus}, can multiply on it, and the associated irritation triggers a defensive discharge as the body tries to expel the object. In a small number of cases, certain strains of \textit{Staphylococcus aureus} produce toxins that enter the bloodstream and cause toxic shock syndrome, a rare but potentially life-threatening illness. The risk of TSS rises the longer a tampon is left in place and with high-absorbency tampons. The discharge caused by a retained tampon is not an infection of the vagina itself at first, but over time the bacterial overgrowth can lead to inflammation and, in rare cases, ascending pelvic infection.

\section*{Clinical Features}
A retained tampon may produce no symptoms at all for days, or even weeks, which is why it is so often discovered by chance. When symptoms do develop, they include:

\begin{itemize}
    \item \textbf{Foul-smelling, often brownish or blood-stained vaginal discharge:} the most common and characteristic symptom, resulting from bacterial overgrowth and tissue irritation
    \item \textbf{Spotting or light bleeding:} may occur even when it is not the time of the period
    \item \textbf{Vaginal itching or soreness:} from local irritation
    \item \textbf{A sensation of something inside the vagina:} some women can feel the tampon or the string
    \item \textbf{Pain during intercourse} or discomfort in the pelvis
    \item \textbf{Offensive odour noticeable to the woman or her partner,} which may be mistaken for poor hygiene
    \item \textbf{In rare cases, symptoms of toxic shock syndrome:} sudden high fever, vomiting, diarrhoea, muscle aches, a sunburn-like rash, dizziness, and confusion
\end{itemize}

Any woman who develops an unexplained vaginal discharge with a bad smell should consider the possibility of a forgotten tampon, especially if she uses tampons.

\section*{Differential Diagnosis}
The symptoms of a retained tampon closely resemble other causes of vaginal discharge and odour, and the list should be considered so that the true cause is not missed.

\begin{itemize}
    \item \textbf{Bacterial vaginosis:} thin grey discharge with a fishy smell, caused by a bacterial imbalance rather than a foreign body
    \item \textbf{Vulvovaginal candidiasis (thrush):} thick white curd-like discharge with intense itching
    \item \textbf{Trichomoniasis or other sexually transmitted infections:} frothy, discoloured discharge, often with itch and soreness
    \item \textbf{Other retained objects:} condoms, contraceptive diaphragms, menstrual cups, or pieces of contraceptive sponge
    \item \textbf{Cervicitis or pelvic inflammatory disease:} discharge with pelvic pain and fever
    \item \textbf{Endometritis or a missed miscarriage:} bleeding and discharge related to the uterus
    \item \textbf{Physiological discharge:} normal mucus changes around ovulation
\end{itemize}

Because the smell and discharge of a retained tampon are so similar to those of an infection, the tampon is often overlooked until it is specifically looked for.

\section*{When to Seek Medical Care}
A woman who realises a tampon is stuck should try to remove it herself; if this fails, or if she is uncomfortable or unable, she should see a general practitioner, sexual health clinic, or pharmacist promptly. It is not an emergency if she feels well, but it should not be left for long. Medical review is appropriate if:

\begin{itemize}
    \item She cannot remove the tampon despite trying
    \item There is an offensive discharge or it has been in place for more than a few days
    \item She is unsure whether a tampon is still present
    \item She has symptoms of a pelvic infection, such as pelvic pain or fever
    \item She is pregnant
\end{itemize}

\textbf{Contact your local emergency services} or go to an emergency department immediately for symptoms of toxic shock syndrome:
\begin{itemize}
    \item Sudden high fever
    \item Vomiting or diarrhoea
    \item A sunburn-like red rash
    \item Dizziness, confusion, or fainting
    \item Muscle aches and weakness
\end{itemize}

Toxic shock syndrome can progress quickly and is a medical emergency even though it is very rare.

\section*{Diagnosis and Investigations}
The diagnosis of a retained tampon is usually made instantly on a speculum examination, which allows the clinician to see the tampon, its string, or its absence. The examination also reveals the cause of any discharge and allows a thorough check that nothing else, such as a second tampon, remains. In most cases no tests are required. If the discharge is severe, offensive, or the tissues look inflamed, a swab may be taken to identify any secondary infection, and occasionally a course of antibiotics is prescribed. If there is any doubt about the presence of a foreign body higher in the vagina, or if a string cannot be seen and the woman is certain a tampon is present, an ultrasound can occasionally be helpful, but this is rarely needed. Once the tampon has been removed, symptoms normally settle quickly without further investigation.

\section*{Management}
Management of a retained tampon is straightforward and starts with removal.

\begin{itemize}
    \item \textbf{Self-removal:} the woman should wash her hands, squat or sit, and gently feel for the string or the tampon with a clean finger. If the string is accessible, a gentle pull usually removes the tampon easily
    \item \textbf{Removal in clinic:} a clinician removes the tampon with forceps during a speculum examination; this takes seconds and is painless for most women
    \item \textbf{Reassurance:} once removed, the overwhelming majority of women need no further treatment
    \item \textbf{Antibiotics:} if there are signs of secondary infection, or if the tampon has been retained for a long time, a short course of antibiotics may be prescribed
    \item \textbf{Treatment of toxic shock syndrome:} if TSS develops, it is treated in hospital with removal of the tampon, intravenous fluids, antibiotics, and supportive care
    \item \textbf{Follow-up:} if symptoms such as discharge persist after removal, a review is arranged; persistent symptoms suggest an alternative cause such as BV or thrush that needs its own treatment
\end{itemize}

It is important to remember that nothing more is usually needed than removal. Many women fear serious harm has occurred, but with prompt removal the outcome is excellent.

\section*{Complications}
The longer a tampon stays in place, the more likely complications become, although even long-retained tampons often cause little harm.

\begin{itemize}
    \item \textbf{Malodorous discharge and vaginitis:} by far the most common consequence, from bacterial overgrowth and tissue irritation
    \item \textbf{Toxic shock syndrome (TSS):} the most feared complication, caused by toxin-producing \textit{Staphylococcus aureus}. It is rare but serious, with fever, rash, shock, and multi-organ involvement, and it is directly linked to prolonged tampon retention
    \item \textbf{Secondary bacterial vaginosis:} the disruption of the normal flora can trigger BV
    \item \textbf{Cervical or vaginal abrasions:} from attempts at removal or from the tampon itself
    \item \textbf{Ascending pelvic infection:} very rare, but possible if the tampon remains for weeks and bacteria travel upward
    \item \textbf{Emotional distress and embarrassment:} women often feel ashamed or worried, though the problem is very common and easily resolved
\end{itemize}

\section*{Prognosis and Outcomes}
The prognosis after a retained tampon is excellent. Once the tampon is removed, the discharge and smell usually resolve within a few days, and no lasting harm is done in the vast majority of cases. Even tampons retained for weeks are typically removed without complication. Toxic shock syndrome is the exception: it is a serious illness with a significant mortality rate even with prompt hospital treatment, which is why it must never be ignored. However, because TSS is rare and the warning signs are well publicised, the practical message is reassuring: remember to change and remove tampons, know the TSS symptoms, and seek help early, and the risk is negligible. Fertility and general health are unaffected by a retained tampon that is removed in time.

\section*{Prevention and Self-Care}
\begin{itemize}
    \item Use the lowest absorbency tampon needed for your flow, and change it every four to eight hours
    \item Make a habit of checking for the string each time you use the toilet during your period
    \item Never leave a tampon in overnight for more than eight hours, and consider a pad or menstrual cup for sleep
    \item Set a phone reminder at the end of each period to confirm the last tampon has been removed
    \item Always remove the final tampon of a period, even if bleeding has stopped
    \item Never insert a tampon outside the time of a period for discharge, as it is easy to forget
    \item Do not insert more than one tampon at a time
    \item Wash your hands before insertion and removal
    \item If the string is lost, do not panic: try to remove the tampon gently with clean fingers, and seek help if that fails
    \item Consider alternating tampons with pads or period underwear, which carry no risk of retention
\end{itemize}

\section*{Sources and Further Reading}
The following organisations provide reliable information about tampon safety, retained tampons, and toxic shock syndrome.

\begin{itemize}
    \item NHS. \textit{Toxic shock syndrome: overview and symptoms.} https://www.nhs.uk/
    \item Mayo Clinic. \textit{Menstrual cycle products and toxic shock syndrome information.} https://www.mayoclinic.org/
    \item MSD Manual. \textit{Toxic shock syndrome: consumer overview.} https://www.msdmanuals.com/
    \item MedlinePlus. \textit{Toxic shock syndrome health information.} https://medlineplus.gov/
    \item Australian Department of Health and Aged Care. \textit{Tampon safety information and TSS warnings.} https://www.health.gov.au/
    \item Jean Hailes for Women's Health. \textit{Period products and menstrual health information.} https://www.jeanhailes.org.au/
\end{itemize}
